\chapter[remarques sur la r\`egle Anneal]%
{Quelques remarques sur la r\`egle Anneal}
\mysetref{about-anneal}

L'observation de la r\`egle~\arule{Anneal} nous conduit \`a proposer l'existence
d'une relative invariance spatiotemporelle des trajectoires dans
l'espace des configurations d'espaces pour des configurations initiales
de densit\'e proche de\,\undemi. %
Nous avons mesur\'e la densit\'e et la taille des fronti\`eres des espaces
cellulaires pour des tailles d'espaces de $2^{20}$~cellules donnant
des transitoires de plus de 100000~g\'en\'erations. %
Nous avons d\'ecouvert qu'il existe des cycles limites dont la taille
est de l'ordre de grandeur du rayon de l'espace cellulaire. %
Nous proposons plusieurs exp\'eriences permettant de compl\'eter cette
\'etude. %

\section{Trajectoires et espaces de trajectoires}

\mydumpfull[p]{anneal-3d}%
{}%
{Une vue~3D de la dynamique de la r\`egle~\arule{Anneal}}%
{Cette image sugg\`ere l'int\'er\^et d'une \'etude des textures produites par
  les plongement~3D d'AC~2D dont la dynamique, \`a l'image
  de~\arule{Anneal}, tend \`a produire des fronti\`eres.}

\mydumptwothirdsR[t]{.75}{life-00}%
{Une signature cellulaire}%
{Exemple de signature d'espace cellulaire.}%
{\arule{Life} sur un espace de taille $128^2$, au temps 128, pour un espace
  initial de densit\'e al\'eatoire\,\Sundemi.}

\mydumptwothirdsR[t]{.75}{anneal-00}%
{}%
{Une configuration typique de la r\`egle~\arule{Anneal}}%
{\arule{Anneal} sur un espace de taille $256^2$, au temps 256,
  pour un espace initial de densit\'e al\'eatoire\,\Sundemi.}

\mydumpfull[t]{anneal-03}%
{}%
{Une s\'equence typique de la r\`egle~\arule{Anneal}}%
{559 g\'en\'erations de~\arule{anneal} sur un espace de taille $256^2$,
  pour un espace initial de densit\'e al\'eatoire\,\Sundemi.}

\mydumpfull[t]{anneal-02}%
{}%
{Deux configurations de la r\`egle~\arule{anneal}}%
{anneal sur deux espaces initiaux de densit\'e al\'eatoire\,\Sundemi.}

\mydumpfull[t]{anneal-01}%
{}%
{Les deux configurations de la figure~\protect\mygetdump{anneal-02} non retouch\'ee}%
{Contrairement \`a la figure~\mygetdump{anneal-00}, les configurations
  locales stables ne sont plus visibles sans perte d'information.}

\myplotfull[t]{anneal-01}%
{}%
{\'Evolution de la taille des fronti\`eres de~\arule{anneal}}%
{La taille des fronti\`eres est une approximation par la diff\'erence bit
  \`a bit entre l'espace de g\'en\'eration~$n$ et~$n-1$.}

\myplotfull[t]{anneal-02}%
{}%
{\'Evolution de la population de~\arule{Anneal} lors
  d'un changement du nombre d'\^iles}%
{Cette figure permet de comprendre l'existence des points d'inflexions
  de la courbe d'\'evolution de anneal.}

\myplotfull[p]{anneal-03}%
{}%
{Transition vers la derni\`ere \^ile}%
{Une autre zone de transition sur les m\^emes donn\'ees que celle de la
  figure~\mygetplot{anneal-02}.}

\myplotfull[p]{anneal-04}%
{}%
{Populations normalis\'ees pour 4~espaces initiaux de tailles crois\-santes}%
{R\'eductions int\'egrales normalis\'ees des populations des trajectoires
  associ\'ees \`a 4~espaces initiaux de tailles respectives:~$2^{128}$,
  $2^{256}$, $2^{512}$ et $2^{1024}$. %
  Le plot de gauche pr\'esente le logarithme de la taille des frontieres.}

\myplotfullR[p]{.8}{anneal-05}%
{}%
{Une s\'erie de zooms sur la population et la taille des fronti\`eres}%
{Les quatres plots du haut pr\'esentent quatres zooms successifs sur la
  population et la s\'equence tri\'ee des populations. %
  Les quatres plots du bas zooment sur la taille des fronti\`eres.}

Les objets de notre \'etude sont les automates cellulaires~(AC). %
Les AC sont des syst\`emes dynamiques discrets. %
Le temps, l'espace et les cellules sont des entit\'es discr\`etes. %
Une fonction de transition~(r\`egle) est appliqu\'ee de mani\`ere synchrone
sur chacune des cellules de l'espace cellulaire. %
Cette fonction re\c{c}oit en argument une cellule et un ensemble de
cellules voisines~(constituant le voisinage pour une r\`egle) %
au temps~$t_n$. %
Elle calcule le nouvel \'etat de la cellule~(au temps~$t_{n+1}$). %

La s\'equence des~$E_i$ espaces g\'en\'er\'ee par une r\`egle~$R$ sur un espace
initial~$E_0$ consid\'er\'e comme une trajectoire dans l'espace
des~$E_{s^d}$~(ou~$s$ est la taille d'un espace cellulaire fini, %
et~$d$ le nombre de dimension cart\'esienne de cet espace) \'etats d'espaces
constitue un des objets \'el\'ementaires de notre \'etude des~AC\@. %
La r\`egle de transition consid\'er\'ee comme une table en fournit un autre. %
La taille de cette table est~$k^{2dr+1}$~(pour une g\'en\'eralisation du
voisinage de Moore), %
$k$ \'etant le nombre d'\'etats possibles pour une cellule~%
(en pratique~$k=2^p$ ou $p$~est le nombre de plans de bits utile de
l'espace), $d$~le nombre de dimensions de l'espace et $r$~le rayon du
voisinage. %
Le passage en param\`etre de~$E_0$ permet de consid\'erer l'ensemble des
trajectoires. %
Le passage en param\`etre de~$R$ introduit l'espace des r\`egles. %
Le nombre de r\`egles est:
%%
\begin{equation}
  Rn_{kdr} = k^{k^{2dr+1}}. %
  \label{rule-count}
\end{equation}
%%
Nous proposons l'\'etude de ces objets par le calcul des trajectoires,
pour des familles \'echantillons sur les param\`etres~$E_0$ et~$R$. %

\section{Vari\'et\'e des objets cellulaires}

Une premi\`ere question porte sur la nature et la vari\'et\'e des~$E_i$
espaces. %
L'espace initial~$E_0$ peut \^etre quelconque et sera en
g\'en\'eral\footnote{%
  Statistiquement au moins.} %
un Jardin d'Eden pour la r\`egle~$R$. %
Empiriquement, les~$E_i$ espaces poss\`edent une signature qui
les rend ais\'ement reconnaissable par l'observateur. %
La figure~\mygetdump{life-00} par exemple sera associ\'ee \`a~\arule{Life}
sans h\'esitation. %
Il convient cependant de remarquer qu'une telle appr\'eciation ne
renseigne que sur la vraisemblance de l'appartenance d'un espace \`a une
trajectoire. %
Dans le cas le plus g\'en\'eral la question de l'appartenance d'un espace
\`a la trajectoire d'une r\`egle~(hors Jardin d'Eden) n'a pas de r\'eponse
en temps polyn\^omial. %
En effet si la r\`egle est universelle~(ce qui est par exemple la cas
de~\arule{Life}) la d\'ecision peut se ramener \`a l'\'enum\'eration des
trajectoires. %

On reporte la question de la vari\'et\'e sur les trajectoires. %
Les s\'equences peuvent \^etre consid\'er\'es comme un espace de dimension
sup\'erieure et on recherche alors les formes
associ\'ees \`a une r\`egle. %
La localit\'e du voisinage introduit une vitesse limite de propagation
de l'information d'une cellule \`a une autre~(une vitesse de la lumi\`ere
cellulaire \'egale au rayon du voisinage). %
L'axe du temps est un axe d'observation privil\'egi\'e, et les trajectoires
peuvent \^etre r\'eduites par contraction de l'espace sur l'axe du temps. %
On utilise une fonction de r\'eduction du nombre de bits d'un espace,
qui joue le r\^ole d'une mesure effectu\'ee sur un syst\`eme physique, %
pour obtenir une s\'equence num\'erique temporelle. %
Le probl\`eme de la signature temporelle d'une r\`egle peut alors \^etre
pos\'e en termes statistiques. %
La mesure de la densit\'e d'un espace~%
(dans le cas le plus g\'en\'eral, si $k > 2$, on peut consid\'erer un
histogramme de $2^p$ \'etat ou $p$ densit\'es) %
est la mesure la plus simple permettant d'obtenir une trace de la
trajectoire d'une r\`egle appliqu\'ee sur une espace initial. %
Le nombre des \'etats d'espaces:
%%
\begin{equation}
  En_{ksd} = k^{s^d}
  \label{space-count}
\end{equation}
%%
impose la r\'eduction des espace initiaux \`a un sous-ensemble. %
Celui-ci peut \^etre fourni par la fonction inverse d'une fonction de
r\'eduction de l'information d'un espace, utilis\'ee comme fonction
d'\'echantillonnage de l'espace des configurations initiales. %
L'\'etude de l'\'evolution de la densit\'e par exemple, est fonction d'un
param\`etre de densit\'e initiale, qui est l'inverse de la fonction de
r\'eduction utilis\'ee, et qui permet d'effectuer un balayage statistique
des \'etats initiaux possibles associ\'e \`a une distribution arbitraire. %
Enfin l'extension de la question de la vari\'et\'e \`a l'espace des r\`egles
(qui combine les difficult\'es combinatoires des \'equation~\ref{rule-count}
et~\ref{space-count}) doit \'egalement s'appuyer sur l'utilisation de
fonction de r\'eduction de l'information associ\'ee \`a une r\`egle, permettant
le balayage des espaces de r\`egles. %

Nous pensons n\'eanmoins qu'il serait utile de disposer d'un
d\'eveloppement complet du calcul pour des espaces cellulaires de
tailles r\'eduites, et des espaces de r\`egles contraintes\footnote{%
  Les r\`egles additives, dont~\arule{Life} et~\arule{Anneal} sont des
  repr\'esentants exemplaires, semble un bon point de d\'epart.}, %
en vue de l'obtention d'indices sur la nature des sym\'etrie, et des
tendances structurelles de l'espace des param\`etres permettant d'orienter
le calcul sur les r\'egions sensibles des espace de tailles
combinatoires. %

\section[r\'esultat exp\'erimentaux]%
{Quelques r\'esultats exp\'erimentaux \`a propos de Anneal}

La r\`egle~\arule{Anneal} est une r\`egle additive\footnote{%
  La r\`egle est une fonction de la somme du voisinage. %
  \mygetrule[phrase]{anneal}} %
dont la table secondaire est:
\begin{verbatim}
0000101111
\end{verbatim}
La figure~\mygetdump{anneal-00} montre une configuration typique
de cette r\`egle pour un espace initial de densit\'e\,\undemi, %
la figure~\mygetdump{anneal-03} une s\'equence typique de
cette r\`egle pour un espace initial de densit\'e\,\undemi %
et la figure~\mygetdump{anneal-02} montre deux configurations de
\arule{Anneal} pour deux espaces initiaux de densit\'e\,\undemi. %

Une propri\'et\'e int\'eressante de cette r\`egle est que des grandeurs
caract\'eristiques des espaces g\'en\'er\'es, comme le nombres d'\^iles ou la
tailles des fronti\`eres, sont ind\'ependantes de l'\'echelle d'observation
choisie, \`a condition de coupler le changement d'\'echelle sur l'espace
\`a un changement d'\'echelle sur le temps. %
La taille et l'\^age des espaces de la figure~\mygetdump{anneal-02}
ne peuvent \^etre d\'etermin\'es simplement. %
En fait les espaces de la figure~\mygetdump{anneal-02} ont \'et\'e
retouch\'es pour faire dispara\^itre les configurations locales stables de
quelque cellules, qui permettent de d\'eterminer l'\'echelle d'espace. %
On admettra que cet effet tend \`a dispara\^itre avec l'augmentation de
tailles des espaces, et l'apparition d'un seuil au del\`a duquel le
groupe isol\'e de quelques cellules n'est plus perceptible~(par l'{\oe}il
ou l'outil de mesure). %
L'espace \`a gauche de la figure~\mygetdump{anneal-02} est de
taille~$1024^2$, son \^age 4184~g\'en\'erations, %
l'espace droit est de taille~$360^2$, son \^age 600~g\'en\'erations.

\noindent%
La figure~\mygetdump{anneal-01} montre les deux configurations de
la figure~\mygetdump{anneal-02} non retouch\'ee. %
On peut remarquer que, d\'ej\`a \`a cette \'echelle, les configurations
locales stables ne sont plus visibles sans perte d'information,
contrairement \`a l'\'echelle de la figure~\mygetdump{anneal-00}, ou de
telles configurations sont visibles. %

Toutes les trajectoires pr\'esent\'ees ici font appara\^itre une diminution de
la densit\'e\footnote{%
  Appel\'ee population sur les plots.}. %
Cette densit\'e est dans le cas de~\arule{Anneal} une quantit\'e
inversible\footnote{%
  Au sens d'une inversion photo-n\'egative, d'un inverseur bit \`a bit appliqu\'e
  sur un espace cellulaire avant l'op\'eration de comptage.}, %
La valeur du param\`etre $\lambda$\footnote{%
  Le param\`etre $\lambda$ est d\'efini comme le pourcentage d'entr\'ees de
  la table de transition dont la valeur est diff\'erente de l'\'etat z\'ero. %
  $\lambda = \frac{k^{dr+1}-\#0}{k^{dr+1}}$}~\cite{Li90a,Li90b}
de~\arule{Anneal} est exactement\,\undemi. %
Bien que nous ayons trouv\'e jusqu'\`a pr\'esent plus de trajectoires
montrant une diminution de la densit\'e des bits \`a~1, nous pensons qu'un
meilleur contr\^ole sur l'\'ecart de la configuration initiale \`a la
valeur\,\undemi~$s^d$ permettrait de v\'erifier le ratio\,\undemi pour
les diminutions de densit\'es. %

Le plot droit de la figure~\mygetplot{anneal-05} montre les
densit\'es des trajectoires pour 4~espaces initiaux de tailles
$2^{128}$, $2^{256}$, $2^{512}$ et $2^{1024}$. %
La premi\`ere atteint un point fixe en 3563~g\'en\'erations, les autres
atteignent un cycle limite de taille $5$, $768$ et~$24$ apr\`es~$6008$,
$29442$ et $107869$~g\'en\'erations. %
Les cycles limites des trajectoires~2 et~4 sont inf\'erieurs d'un ordre
de grandeur au rayon de l'espace, alors que celui de la trajectoire~3
est du m\^eme ordre de grandeur que le rayon de l'espace. %
De plus, il est visible sur le plot, que la densit\'e finale du
transitoire est tr\`es importante pour la trajectoire~3 contrairement
aux autres, pour lesquels, tant la taille du cycle que la densit\'e
moyenne de celui-ci peuvent \^etre consid\'er\'ees comme une valeur approch\'ee
d'un des deux points fixes majoritaires de~\arule{Anneal}. %
Les densit\'es initiales exactes sont pour chaque trajectoire
respectivement:~%
$8172/128^2 = 0.498779296875$, %
$32673/256^2 = 0.4985504150390625$, %
$130960/512^2 = 0.49957275390625$, %
$523619/1024^2 = 0.49936199188232421875$, %
L'\'ecart \`a\,\undemi \'etant toujours n\'egatif peut expliquer le destin
vers z\'ero de toutes ces configurations. %
La question de la sensibilit\'e de~\arule{Anneal} aux \'ecarts de la
densit\'e\,\undemi, et en particulier du rapport entre cette densit\'e et
le nombre de trajectoire n'atteignant pas un quasi point fixe, reste
pour nous ouverte. %

Une question sugg\'er\'ee par ces exp\'eriences est celle de la taille des
transitoires de~\arule{anneal} au regard de la taille moyenne des
transitoires de la famille de r\`egle associ\'ee \`a un voisinage. %
Le probl\`eme de la fabrication d'une r\`egle de taille de
transitoire donn\'e~(\'eventuellement maximum), ne poss\`ede pas de r\'eponse
simple. %
Nous proposons de consid\'erer la r\`egle~\arule{anneal} comme prototype
pour l'\'etude de cette question. %

\noindent%
La figure~\mygetplot{anneal-01} montre l'\'evolution de la taille
des fronti\`eres d'une trajectoire de~\arule{Anneal} pour un espace de
taille~$1024^2$ et de densit\'e initiale\,\undemi. %
La figure~\mygetplot{anneal-02} montre sur les m\^emes donn\'ees que
celle de la figure~\mygetplot{anneal-01}, une transition sur le
nombre d'\^iles et les variations de la population associ\'ee. %

%%
%% Local Variables:
%% mode: auto-fill
%% iso-accents-minor-mode: nil
%% TeX-master: "top"
%% TeX-command-default: "mytex"
%% Local IspellParsing: "latex-mode ~latin1"
%% Local IspellPersDict: "~/.ispell_lisp"
%% LocalWords: "pt htb dump eps IspellParsing qwf dummy Moore width latex-mode"
%% LocalWords: "IspellPersDict ispell lisp LocalWords End francais english babel"
%% LocalWords: "d'espaces d'eden eden anneal paper dvips afterpage epsfig email"
%% LocalWords: "Thierry Delamare thy tao.univ-paris.fr transitoires height plain"
%% LocalWords: "not DENIS univ fr about-anneal d'AC d'Eden photo-n\'egative"
%% LocalWords: "santes zooment"
%% End:
